% Generated draft: requires mathematical and source review.
\documentclass[11pt]{article}
\usepackage[margin=25mm]{geometry}
\usepackage{fontspec}
\setmainfont{Latin Modern Roman}
\usepackage{amsmath,amssymb,mathtools,unicode-math}
\setmathfont{Latin Modern Math}
\usepackage{enumitem,needspace,xurl,hyperref}
\setlength{\parindent}{0pt}
\setlength{\parskip}{5pt}
\setlength{\emergencystretch}{3em}
\newcommand{\N}{\mathbb{N}}
\newcommand{\Z}{\mathbb{Z}}
\newcommand{\Q}{\mathbb{Q}}
\newcommand{\R}{\mathbb{R}}
\newcommand{\C}{\mathbb{C}}
\newcommand{\F}{\mathbb{F}}
\newcommand{\A}{\mathbb{A}}
\newcommand{\PP}{\mathbb{P}}
\newcommand{\E}{\mathbb{E}}
\newcommand{\eps}{\varepsilon}
\newcommand{\dd}{\,\mathrm{d}}
\begin{document}
\section*{UnsolvedMath ID: 700881}
\textbf{Source problem code:} catalogue entry
\subsection{Definitions and mathematical statement}
Let \(\Omega\) be an infinite set and let \(G\leq\operatorname{Sym}(\Omega)\) be a permutation group, acting faithfully on \(\Omega\). Define \(G^{(0)}=G\) and \(G^{(i+1)}=[G^{(i)},G^{(i)}]\), where \([H,H]\) denotes the subgroup generated by all commutators of elements of \(H\). The group \(G\) is soluble if \(G^{(d)}=\{1\}\) for some finite integer \(d\geq 0\).

For an integer \(k\geq 1\), put
\[
\Omega^{\langle k\rangle}=\{(\omega_1,\ldots,\omega_k)\in\Omega^k:\omega_i\neq\omega_j\text{ whenever }i\neq j\}.
\]
The diagonal action is
\[
g\cdot(\omega_1,\ldots,\omega_k)=(g\omega_1,\ldots,g\omega_k).
\]
Write \(G\backslash\Omega^{\langle k\rangle}\) for its set of orbits.

\textbf{Problem (P. M. Neumann, 1992).} Does there exist an infinite set \(\Omega\) and a soluble subgroup \(G\leq\operatorname{Sym}(\Omega)\) such that
\[
\bigl|G\backslash\Omega^{\langle 3\rangle}\bigr|<\infty?
\]
\textbf{Supplied conjecture.} No such group exists; equivalently, every soluble permutation group on an infinite set has infinitely many orbits on triples. No countability, transitivity, or fixed bound on the derived length is assumed.

\textbf{Qualification concerning triples.} The supplied wording does not specify whether triples are ordered, unordered, or allowed to have repeated entries. We use ordered triples of distinct points. On an infinite set, finiteness of the number of orbits is equivalent for these three conventions, so this choice does not change the existence question.

\textbf{Supplied comment of 2001.} The record reports that D. Macpherson proved that every infinite soluble permutation group has infinitely many orbits on quadruples, citing the 1997 reference below. This is retained as a source-reported result, without independent verification; it does not settle the stated triples question.
\subsection{Short English statement}
Can a soluble group act faithfully on an infinite set with only finitely many orbits on triples of points? P. M. Neumann conjectured that the answer is no. The supplied record labels the problem open and includes a 2001 comment reporting an analogous result for quadruples.
\subsection{Sources}
\begin{itemize}
\item Supplied local UnsolvedMath registry record: permanent ID \(700881\), exact source code \texttt{catalogue entry}, title \emph{Kourovka Notebook Problem 12.63}. The record was created and updated on \texttt{2026-10-16T14:00:00Z}; its supplied status is \texttt{open}. Upstream revision: \texttt{d371ee917ac9647fbb06cd6fc529deec6f3dda82}.
\item \emph{The Kourovka Notebook (Unsolved Problems in Group Theory)}, original problem 12.63, 12th issue (1992), proposed by P. M. Neumann. The supplied background identifies the 21st issue edition (2026), \texttt{arXiv:1401.0300v48}; the collection metadata names E. I. Khukhro and V. D. Mazurov as editors.
\item D. Macpherson, contribution in \emph{Advances in algebra and model theory}, Algebra Logic Appl., volume 9, Gordon \& Breach, Amsterdam, 1997, pp. 87--92. This is the reference supplied by the comment of 2001 for the assertion about quadruples; no contribution title is supplied.
\end{itemize}
\end{document}
